
% MemSys Spec
\label{sec:memsysfunc01}

This section specifies the SystemVerilog module decomposition of the memory subsystem. Each module is described with its complete port list (signal, direction, width, connected to/from, and purpose) and its internal algorithm. The module names follow the project-wide \texttt{as} prefix convention.

\paragraph{Module Hierarchy:}
\label{subsubsec:modHierarchy}

Table~\ref{tab:asModHierarchya} lists all modules.
\texttt{asMemTop} is the sole instantiation point; the CPU core and the
QSPI controller are outside this boundary.

\begin{table}[H]
\caption{Memory Subsystem -- SystemVerilog Module Overview}
\label{tab:asModHierarchya}
\centering
\begin{tabularx}{\textwidth}{|l |l |X|}
  \hline
  Module & Type & Purpose \\
  \hline
  \hline
  \texttt{asMemTop}      & top wrapper   & Instantiates all sub-modules; routes CPU, SRAM and AXI4 interfaces \\
  \hline
  \texttt{asICacheCtrl}  & RTL (FSM)     & Instruction cache controller: tag compare, PLRU, AXI4 master (AR/R) \\
  \hline
  \texttt{asDCacheCtrl}  & RTL (FSM)     & Data cache controller: tag compare, PLRU, bypass path, AXI4 master (AR/R) \\
  \hline
  \texttt{asICacheBlock} & SRAM wrapper  & Holds I-Tag SRAM (96\,bit$\times$32) and I-Data SRAM (256\,bit$\times$128) \\
  \hline
  \texttt{asDCacheBlock} & SRAM wrapper  & Holds D-Tag SRAM (96\,bit$\times$32) and D-Data SRAM (256\,bit$\times$128) \\
  \hline
  \texttt{asSpBlock}     & SRAM wrapper  & Holds Scratchpad SRAM (64\,bit$\times$1024) \\
  \hline
  \texttt{asMemArb}      & RTL (mux)     & Fixed-priority AXI4 arbiter: D-Cache\,$>$\,I-Cache, single output port \\
  \hline
  \texttt{asSpramModel}  & sim. model    & Behavioural 1-cycle synchronous SRAM; replaces X-Fab SPRAM macro in simulation \\
  \hline
\end{tabularx}
\end{table}

The three SRAM wrapper modules (\texttt{asICacheBlock}, \texttt{asDCacheBlock},
\texttt{asSpBlock}) each instantiate one or two X-Fab XO035 SPRAM macros.
In simulation the macros are replaced by \texttt{asSpramModel} via a
synthesis guard (\texttt{`ifdef SIMULATION}).

%%%%%%%%%%%%%%%%%%%%%%%%%%%%%%%%%
% System Integration of MemSys
%%%%%%%%%%%%%%%%%%%%%%%%%%%%%%%%%
\subsection{System Integration of MemSys0}
%%%%%%%%%%%%%%%%%%%%%%%%%%%%%%%%%
% System Integration of MemSys0: History
%%%%%%%%%%%%%%%%%%%%%%%%%%%%%%%%%
\subsubsection{History}
Regulations for author history:
\begin{itemize}
  \item Only visible internally (by conditioning of the table, no other conditions needed)
\end{itemize}

\begin{table}[H]
\caption{Authors}
\label{tab:memsysauth01a}
\centering
\begin{tabularx}{\textwidth}{|X |X |X |}
  \hline
  Author & Department & From \\
  \hline
  \hline
  Andreas Siggelkow & Fac. E & February 2026 \\
  \hline
  && \\
  \hline
\end{tabularx}
\end{table}

Regulations for history table:
\begin{itemize}
  \item Last version on top.
  \item Dont use links in the history table. These may cause unexpected history changes in future versions.
  \item Name (or short sign) should be visible only internally.
  \item Create a new version after each (also limited) release of the document. Versioning does not depend on the overall document version.
  \item Use the form x.y for the version. x for the master versions, y for the project specific changes and updates.
\end{itemize}

\begin{table}[H]
\caption{History}
\label{tab:memsyshistory01a}
\centering
\begin{tabularx}{\textwidth}{|X |X |X |}
  \hline
  Date & Name & Content Changes \\
  \hline
  \multicolumn{3}{|l|}{Version of this document: V0.1} \\
  \hline
  2026-05-14 & Andreas Siggelkow & First draft (AS, Claude Code) \\
  \hline
\end{tabularx}
\end{table}

%%%%%%%%%%%%%%%%%%%%%%%%%%%%%%%%%
% System Integration of MemSys0: Functional Configuration of QSPI0
%%%%%%%%%%%%%%%%%%%%%%%%%%%%%%%%%
\subsubsection{Functional Configuration of MemSys0}
\textcolor{red}{This section should describe the product specific instantiation of this module’s feature set (which may be a subset of the full feature set). It is intended also for use in the product overview section of the target spec/summary target spec.}

The standard features of a MemSys peripheral are described in its functional description in Section \ref{subsec:memsysfunc01}. Generic features, that can be decided for each specific instantiation of the peripheral are listed and completed for MemSys0 in Table \ref{tab:asMemSys01a}.
\begin{table}[H]
\caption{Memory System Feature Set}
\label{tab:asMemSys01a}
\centering
\begin{tabularx}{\textwidth}{|l |l |X|}
  \hline
   & Generic MemSys Feature Set & Details \\
  \hline
  \hline
  yes & Separate Kernel Clock      & 125\;MHz (TBD) \\
  \hline
\end{tabularx}
\end{table}

%%%%%%%%%%%%%%%%%%%%%%%%%%%%%%%%%
% System Integration of MemSys0: HW Interfaces
%%%%%%%%%%%%%%%%%%%%%%%%%%%%%%%%%
\subsubsection{HW Interfaces}
The following table \ref{tab:asMemSys02} lists all hardware interfaces which are relevant for the regular operation of the device. It is not a detailed list of all signals and does not include implementation specific informations.

\textcolor{red}{The names used below show the functional connectivity. There is no guarantee that they are electrically compatible (e.g. voltage levels, clock domains etc. may be different). A full list is part of the design spec.}
\begin{table}[H]
\caption{HW Interfaces of Module MemSys0}
\label{tab:asMemSys02}
\centering
\begin{tabularx}{\textwidth}{|l |l |X|}
  \hline
   & Module internal name & Chip level name \\
  \hline
  \hline
  Supply &  &  \\
  \hline
  Clock & clk\_i & clk\_i \\
  \hline
  I-Cache-Bus  & I-Cache-Bus slave & - \\
  \hline
  D-Cache-Bus  & D-Cache-Bus slave & - \\
  \hline
  Interrupt & None &  - \\
  \hline
  DMA & None & - \\
  \hline
  External Signals & None & - \\
  \hline
\end{tabularx}
\end{table}

%%%%%%%%%%%%%%%%%%%%%%%%%%%%%%%%%
% System Integration of MemSys0: Bus Interface
%%%%%%%%%%%%%%%%%%%%%%%%%%%%%%%%%
\subsubsection{Bus Interface}
The memory system  exposes three independent bus interfaces: AXI, I-Mem-IF, D-Mem-IF.


%%%%%%%%%%%%%%%%%%%%%%%%%%%%%%%%%
% System Integration of MemSys0: Registers
%%%%%%%%%%%%%%%%%%%%%%%%%%%%%%%%%
\subsubsection{Registers}
No registers.

%%%%%%%%%%%%%%%%%%%%%%%%%%%%%%%%%
% System Integration of MemSys0: Related Sections of this Spec
%%%%%%%%%%%%%%%%%%%%%%%%%%%%%%%%%
\subsubsection{Related Sections of this Spec}
\begin{itemize}
  \item MemSys functional description starting with Section \ref{subsec:memsysfunc01}.
  %\item Interrupt and DMA Concept Section ``Architecture Concepts''
  \item Peripheral Architecture Concept Section ``Architecture Concepts''
  \begin{itemize}
    \item Slave WB (Wishbone) BPI Section ``Architecture Concepts''
    \item AXI4 Cache Bus Section ``Architecture Concepts'' (Bus Concept)
    \item Peripheral ID Registers Section ``Architecture Concepts''
    \item Peripheral Clocking Concept Section ``Architecture Concepts''
    \item Clock Control Registers Section ``Architecture Concepts''
    %\item Horizontal Interrupt and DMA Register Structure Section ``Architecture Concepts''
    %\item Basic Interrupt and DMA Register Set Section ``Architecture Concepts''
    \item Interrupt Source Combining Register Set Section ``Architecture Concepts''
    \item Synchronization Block Section ``Architecture Concepts''
    \item Standardized Flip-Flop based FIFO Section ``Architecture Concepts''
    \item TX-FIFO Registers Section ``Architecture Concepts''
    \item RX-FIFO Registers Section ``Architecture Concepts''
  \end{itemize}
\end{itemize}

%%%%%%%%%%%%%%%%%%%%%%%%%%%%%%%%%
% History of MemSys
%%%%%%%%%%%%%%%%%%%%%%%%%%%%%%%%%
\subsection{History of MemSys}
\begin{table}[H]
\caption{History}
\label{tab:memsyshistory02}
%\begin{table}[ht]
\centering
\begin{tabularx}{\textwidth}{|X |X |X |}
  \hline
  \multicolumn{3}{|l|}{History} \\
  \hline
  \multicolumn{1}{|l}{Target Spec.} & \multicolumn{1}{l}{Current version : 1.0, 2026-02-17} & \multicolumn{1}{l|}{Author: A. Siggelkow}\\
  \multicolumn{1}{|l}{ } & \multicolumn{1}{l}{Previous version : 0.0, 2020-10-12} & \multicolumn{1}{l|}{ }\\
  \hline
  \hline
  Paragraph (in previous version) & Paragraph (in current version) / date & Subjects (major changes since last revision) \\
  \hline
  \hline
  \multicolumn{3}{|l|}{Current Chapter Version: C0.1; Date: 2026-02-17} \\
  \hline
  \ref{sec:memsysfunc01} & \ref{sec:memsysfunc01} / 2026-05-14 & First draft (AS, Claude Code) \\
  \hline
\end{tabularx}
\end{table}

%%%%%%%%%%%%%%%%%%%%%%%%%%%%%%%%%
% Functional Overview
%%%%%%%%%%%%%%%%%%%%%%%%%%%%%%%%%
\subsection{Functional Overview}
\label{subsec:memsysfunc01}

\paragraph{Purpose:}
Top-level wrapper that wires the CPU interfaces, the SRAM blocks, the cache controllers, and the memory arbiter together. It presents three external interfaces: 
\begin{itemize}
  \item the CPU instruction bus, 
  \item the CPU data bus, and 
  \item the AXI4 master port toward the QSPI controller.
\end{itemize}
No combinatorial or sequential logic resides in this module beyond wire connections and interface bindings.

\paragraph{Internal algorithm:}
\texttt{asMemTop} contains no logic.
It instantiates exactly the seven modules listed in
Table~\ref{tab:asModHierarchy} (excluding \texttt{asSpramModel}, which
lives inside the SRAM wrappers) and connects them via the SV interfaces
\texttt{as\_icache\_if}, \texttt{as\_dcache\_if}, and \texttt{as\_axi4\_if}.
The SRAM wrappers are wired to their respective controllers through
dedicated flat port groups (see Sections~\ref{subsubsec:asICacheBlock}
through~\ref{subsubsec:asSpBlock}).

%%%%%%%%%%%%%%%%%%%%%%%%%%%%%%%%%
% Structural Overview
%%%%%%%%%%%%%%%%%%%%%%%%%%%%%%%%%
\subsection{Structural Overview}
The MemSys  see Chapter "Architectural Concepts - Memory System"

\paragraph{Port list} of the memory top-level:

\begin{table}[H]
\caption{\texttt{asMemTop} Port List (1)}
\label{tab:asMemTopPorts}
\centering
\begin{tabularx}{\textwidth}{|l |l |l |l |X|}
  \hline
  Signal & Dir & Width & Connected to & Description \\
  \hline
  \hline
  \multicolumn{5}{|l|}{\textit{Clock and Reset}} \\
  \hline
  \texttt{clk\_i}  & in & 1 & All sub-modules & System clock (synchronous rising edge) \\
  \hline
  \texttt{rst\_i}  & in & 1 & All sub-modules & Synchronous active-high reset \\
  \hline
  \multicolumn{5}{|l|}{\textit{CPU Instruction Bus (via \texttt{as\_icache\_if})}} \\
  \hline
  \texttt{ic\_addr\_i}        & in  & 64 & \texttt{asICacheCtrl} & Fetch address from CPU IF stage \\
  \hline
  \texttt{ic\_req\_i}         & in  & 1  & \texttt{asICacheCtrl} & Fetch request valid \\
  \hline
  \texttt{ic\_flush\_i}       & in  & 1  & \texttt{asICacheCtrl} & FENCE.I pulse from CPU \\
  \hline
  \texttt{ic\_stall\_o}       & out & 1  & CPU IF stage          & Cache miss / flush in progress \\
  \hline
  \texttt{ic\_flush\_done\_o} & out & 1  & CPU IF stage          & Flush complete pulse \\
  \hline
  \texttt{ic\_err\_o}         & out & 1  & CPU IF stage          & AXI4 read error pulse \\
  \hline
  \texttt{ic\_rdata\_o}       & out & 32 & CPU IF stage          & Instruction word (32-bit) \\
  \hline
  \texttt{ic\_rvalid\_o}      & out & 1  & CPU IF stage          & Instruction word valid \\
  \hline
\end{tabularx}
\end{table}

\begin{table}[H]
\caption{\texttt{asMemTop} Port List (2)}
\label{tab:asMemTopPorts2}
\centering
\begin{tabularx}{\textwidth}{|l |l |l |l |X|}
  \hline
  Signal & Dir & Width & Connected to & Description \\
  \hline
  \hline
  \multicolumn{5}{|l|}{\textit{CPU Data Bus (via \texttt{as\_dcache\_if})}} \\
  \hline
  \texttt{dc\_addr\_i}        & in  & 64 & \texttt{asDCacheCtrl} & Load/store address from CPU MEM stage \\
  \hline
  \texttt{dc\_req\_i}         & in  & 1  & \texttt{asDCacheCtrl} & Load/store request valid \\
  \hline
  \texttt{dc\_wr\_i}          & in  & 1  & \texttt{asDCacheCtrl} & 1\,=\,store, 0\,=\,load \\
  \hline
  \texttt{dc\_size\_i}        & in  & 3  & \texttt{asDCacheCtrl} & Transfer size (byte/half/word/double) \\
  \hline
  \texttt{dc\_wdata\_i}       & in  & 64 & \texttt{asDCacheCtrl} & Store data from CPU \\
  \hline
  \texttt{dc\_wstrb\_i}       & in  & 8  & \texttt{asDCacheCtrl} & Byte enables for store \\
  \hline
  \texttt{dc\_flush\_i}       & in  & 1  & \texttt{asDCacheCtrl} & FENCE.I pulse from CPU \\
  \hline
  \texttt{dc\_stall\_o}       & out & 1  & CPU MEM stage         & Cache miss / bypass / flush in progress \\
  \hline
  \texttt{dc\_flush\_done\_o} & out & 1  & CPU MEM stage         & Flush complete pulse \\
  \hline
  \texttt{dc\_err\_o}         & out & 1  & CPU MEM stage         & AXI4 read error pulse \\
  \hline
  \texttt{dc\_rdata\_o}       & out & 64 & CPU MEM stage         & Load data (64-bit) \\
  \hline
  \texttt{dc\_rvalid\_o}      & out & 1  & CPU MEM stage         & Load data valid \\
  \hline
  \multicolumn{5}{|l|}{\textit{AXI4 Master Port toward QSPI Controller (via \texttt{as\_axi4\_if})}} \\
  \hline
  \texttt{m\_axi4}            & out & -- & \texttt{asQspiTop}    & AXI4 master bundle (AR, R channels only; see Table~\ref{tab:asAxi4Params}) \\
  \hline
\end{tabularx}
\end{table}


%%%%%%%%%%%%%%%%%%%%%%%%%%%%%%%%%%%%%%%%%%%%%%%%%%%%%%%%%%%
\subsubsection{\texttt{asICacheBlock} -- I-Cache SRAM Wrapper}
\label{subsubsec:asICacheBlock}

\paragraph{Purpose.}
Wraps two X-Fab XO035 SPRAM macros: the I-Tag SRAM (96\,bit\,$\times$\,32)
and the I-Data SRAM (256\,bit\,$\times$\,128).
Translates the logical controller signals into the physical SPRAM port
signals (\texttt{CLK}, \texttt{CEN}, \texttt{WEN}, \texttt{BWEN},
\texttt{A}, \texttt{D}, \texttt{Q}) described in
Section~\ref{subsec:sramMacros}.
In simulation the SPRAM macros are replaced by \texttt{asSpramModel}.

\paragraph{Port list} of the I-cache:


\begin{table}[H]
\caption{\texttt{asICacheBlock} Port List}
\label{tab:asICacheBlockPorts}
\centering
\begin{tabularx}{\textwidth}{|l |l |l |l |X|}
  \hline
  Signal & Dir & Width & Connected to & Description \\
  \hline
  \hline
  \texttt{clk\_i}         & in  & 1   & SPRAM CLK         & System clock \\
  \hline
  \multicolumn{5}{|l|}{\textit{Tag SRAM port (I-Tag SRAM, 96\,bit\,$\times$\,32)}} \\
  \hline
  \texttt{it\_addr\_i}    & in  & 5   & \texttt{asICacheCtrl} & Set index (0--31) \\
  \hline
  \texttt{it\_wdata\_i}   & in  & 96  & \texttt{asICacheCtrl} & Tag write data: \{PLRU[2:0], way3\{valid,tag[21:0]\}, \ldots, way0\{valid,tag[21:0]\}\} \\
  \hline
  \texttt{it\_we\_i}      & in  & 1   & \texttt{asICacheCtrl} & Tag write enable (active-high) \\
  \hline
  \texttt{it\_ce\_i}      & in  & 1   & \texttt{asICacheCtrl} & Tag chip enable (active-high) \\
  \hline
  \texttt{it\_rdata\_o}   & out & 96  & \texttt{asICacheCtrl} & Tag read data; valid one cycle after \texttt{it\_ce\_i} on read \\
  \hline
  \multicolumn{5}{|l|}{\textit{Data SRAM port (I-Data SRAM, 256\,bit\,$\times$\,128)}} \\
  \hline
  \texttt{id\_addr\_i}    & in  & 7   & \texttt{asICacheCtrl} & Address: \{index[4:0], way[1:0]\} \\
  \hline
  \texttt{id\_wdata\_i}   & in  & 64  & \texttt{asICacheCtrl} & One 64-bit beat to write; placed into the lane selected by \texttt{id\_wlane\_i} \\
  \hline
  \texttt{id\_wlane\_i}   & in  & 2   & \texttt{asICacheCtrl} & Beat index (0--3); selects 64-bit lane within 256-bit word \\
  \hline
  \texttt{id\_we\_i}      & in  & 1   & \texttt{asICacheCtrl} & Data write enable (active-high) \\
  \hline
  \texttt{id\_ce\_i}      & in  & 1   & \texttt{asICacheCtrl} & Data chip enable (active-high) \\
  \hline
  \texttt{id\_rdata\_o}   & out & 256 & \texttt{asICacheCtrl} & Full 256-bit cache line; valid one cycle after \texttt{id\_ce\_i} on read \\
  \hline
\end{tabularx}
\end{table}

\paragraph{Internal algorithm.}
On a write (\texttt{it\_we\_i} or \texttt{id\_we\_i} high):
\begin{itemize}
  \item Tag SRAM: \texttt{BWEN} is set to all-zero (write all 96 bits);
    \texttt{WEN} is driven low; \texttt{A} receives \texttt{it\_addr\_i};
    \texttt{D} receives \texttt{it\_wdata\_i}.
  \item Data SRAM: \texttt{id\_wlane\_i} selects one 64-bit lane.
    \texttt{BWEN[255:0]} is constructed as: bits\,$[(lane+1)\cdot64-1\,:\,lane\cdot64]$
    are driven to \texttt{0} (write-enable); all other bits are \texttt{1}
    (mask).
    \texttt{D[63:0]} is replicated four times to fill \texttt{D[255:0]};
    only the selected lane is actually written due to BWEN masking.
\end{itemize}
On a read (\texttt{it\_ce\_i} or \texttt{id\_ce\_i} high, write enables
low): \texttt{WEN} is driven high; \texttt{BWEN} is irrelevant (ignored
on reads by the SPRAM); the read data appears on \texttt{Q} one clock
cycle later.

% -------------------------------------------------------------------------
\subsubsection{\texttt{asDCacheBlock} -- D-Cache SRAM Wrapper}
\label{subsubsec:asDCacheBlock}

\paragraph{Purpose.}
Identical in structure to \texttt{asICacheBlock} but serves the D-Cache
controller.
Wraps D-Tag SRAM (96\,bit\,$\times$\,32) and D-Data SRAM
(256\,bit\,$\times$\,128); same SPRAM macros, same timing, same
simulation guard.

\paragraph{Port list} of the D-cache:

\begin{table}[H]
\caption{\texttt{asDCacheBlock} Port List}
\label{tab:asDCacheBlockPorts}
\centering
\begin{tabularx}{\textwidth}{|l |l |l |l |X|}
  \hline
  Signal & Dir & Width & Connected to & Description \\
  \hline
  \hline
  \texttt{clk\_i}         & in  & 1   & SPRAM CLK             & System clock \\
  \hline
  \multicolumn{5}{|l|}{\textit{Tag SRAM port (D-Tag SRAM, 96\,bit\,$\times$\,32)}} \\
  \hline
  \texttt{dt\_addr\_i}    & in  & 5   & \texttt{asDCacheCtrl} & Set index (0--31) \\
  \hline
  \texttt{dt\_wdata\_i}   & in  & 96  & \texttt{asDCacheCtrl} & Tag write data (same layout as I-Cache; no dirty bit) \\
  \hline
  \texttt{dt\_we\_i}      & in  & 1   & \texttt{asDCacheCtrl} & Tag write enable (active-high) \\
  \hline
  \texttt{dt\_ce\_i}      & in  & 1   & \texttt{asDCacheCtrl} & Tag chip enable (active-high) \\
  \hline
  \texttt{dt\_rdata\_o}   & out & 96  & \texttt{asDCacheCtrl} & Tag read data; valid one cycle after chip enable \\
  \hline
  \multicolumn{5}{|l|}{\textit{Data SRAM port (D-Data SRAM, 256\,bit\,$\times$\,128)}} \\
  \hline
  \texttt{dd\_addr\_i}    & in  & 7   & \texttt{asDCacheCtrl} & Address: \{index[4:0], way[1:0]\} \\
  \hline
  \texttt{dd\_wdata\_i}   & in  & 64  & \texttt{asDCacheCtrl} & One 64-bit AXI4 beat (fill path) \\
  \hline
  \texttt{dd\_wlane\_i}   & in  & 2   & \texttt{asDCacheCtrl} & Beat index (0--3) selecting 64-bit lane \\
  \hline
  \texttt{dd\_we\_i}      & in  & 1   & \texttt{asDCacheCtrl} & Data write enable (active-high) \\
  \hline
  \texttt{dd\_ce\_i}      & in  & 1   & \texttt{asDCacheCtrl} & Data chip enable (active-high) \\
  \hline
  \texttt{dd\_rdata\_o}   & out & 256 & \texttt{asDCacheCtrl} & Full 256-bit cache line \\
  \hline
\end{tabularx}
\end{table}

\paragraph{Internal algorithm.}
Identical to \texttt{asICacheBlock}.
The D-Tag SRAM contains no dirty bit (removed; see
Section~\ref{subsubsec:dcacheAddrPolicy}): the 96-bit tag word layout
is \{PLRU[2:0], way3\{valid[0],tag[21:0]\}, way2\{\ldots\},
way1\{\ldots\}, way0\{valid[0],tag[21:0]\}\}, identical to I-Cache.

% -------------------------------------------------------------------------
\subsubsection{\texttt{asSpBlock} -- Scratchpad SRAM Wrapper}
\label{subsubsec:asSpBlock}

\paragraph{Purpose.}
Wraps a single X-Fab XO035 SPRAM macro configured as
64\,bit\,$\times$\,1024 (8\,KiB default; parameterisable via
\texttt{SP\_DEPTH}).
Provides byte-level write access required for sub-word stores (byte,
halfword, word) to stack, heap, \texttt{.data}, and \texttt{.bss}.

\paragraph{Port list} of the scratch-pad memory:

\begin{table}[H]
\caption{\texttt{asSpBlock} Port List}
\label{tab:asSpBlockPorts}
\centering
\begin{tabularx}{\textwidth}{|l |l |l |l |X|}
  \hline
  Signal & Dir & Width & Connected to & Description \\
  \hline
  \hline
  \texttt{clk\_i}        & in  & 1  & SPRAM CLK             & System clock \\
  \hline
  \texttt{sp\_addr\_i}   & in  & 10 & \texttt{asDCacheCtrl} & Word address within Scratchpad (byte address bits [12:3]) \\
  \hline
  \texttt{sp\_wdata\_i}  & in  & 64 & \texttt{asDCacheCtrl} & Write data (64-bit aligned; byte enables select active bytes) \\
  \hline
  \texttt{sp\_wstrb\_i}  & in  & 8  & \texttt{asDCacheCtrl} & Byte write enables (1\,=\,write byte, 0\,=\,mask); driven from \texttt{dc\_wstrb\_i} \\
  \hline
  \texttt{sp\_we\_i}     & in  & 1  & \texttt{asDCacheCtrl} & Write enable (active-high); 0 = read \\
  \hline
  \texttt{sp\_ce\_i}     & in  & 1  & \texttt{asDCacheCtrl} & Chip enable (active-high) \\
  \hline
  \texttt{sp\_rdata\_o}  & out & 64 & \texttt{asDCacheCtrl} & Read data; valid one cycle after chip enable on a read \\
  \hline
\end{tabularx}
\end{table}

\paragraph{Internal algorithm.}
\begin{itemize}
  \item \textbf{Read} (\texttt{sp\_we\_i} low, \texttt{sp\_ce\_i} high):
    SPRAM \texttt{WEN} driven high; \texttt{CEN} driven low (active-low);
    \texttt{A} receives \texttt{sp\_addr\_i};
    \texttt{sp\_rdata\_o} latches \texttt{Q} in the cycle after access.
  \item \textbf{Write} (\texttt{sp\_we\_i} high, \texttt{sp\_ce\_i} high):
    SPRAM \texttt{WEN} driven low; \texttt{CEN} driven low;
    \texttt{BWEN[63:0]} is formed as the bit-expansion of
    \texttt{\textasciitilde sp\_wstrb\_i}: each byte-enable bit expands to
    eight BWEN bits (0\,=\,write, 1\,=\,mask), enabling sub-byte
    granularity.
    \texttt{D} receives \texttt{sp\_wdata\_i}.
  \item \textbf{Idle}: \texttt{CEN} driven high (chip disabled, SPRAM
    in zero-power standby).
\end{itemize}

% -------------------------------------------------------------------------
\subsubsection{\texttt{asICacheCtrl} -- Instruction Cache Controller}
\label{subsubsec:asICacheCtrl}

\paragraph{Purpose.}
Implements the I-Cache FSM (Table~\ref{tab:asICacheFsm}).
Drives the \texttt{asICacheBlock} SRAM interface, accepts instruction
fetch requests from the CPU via \texttt{as\_icache\_if}, and issues
AXI4 AR/R burst transactions to \texttt{asMemArb} on a cache miss.

\paragraph{Port list} of the I-cache controller:

\begin{table}[H]
\caption{\texttt{asICacheCtrl} Port List (1)}
\label{tab:asICacheCtrlPorts}
\centering
\begin{tabularx}{\textwidth}{|l |l |l |l |X|}
  \hline
  Signal & Dir & Width & Connected to & Description \\
  \hline
  \hline
  \texttt{clk\_i}  & in & 1 & \texttt{asMemTop} & System clock \\
  \hline
  \texttt{rst\_i}  & in & 1 & \texttt{asMemTop} & Synchronous reset; clears FSM state and all tag valid bits \\
  \hline
  \multicolumn{5}{|l|}{\textit{CPU Instruction Bus (\texttt{as\_icache\_if.cache} modport)}} \\
  \hline
  \texttt{ic\_addr\_i}        & in  & 64 & CPU IF stage & Fetch address \\
  \hline
  \texttt{ic\_req\_i}         & in  & 1  & CPU IF stage & Fetch request valid \\
  \hline
  \texttt{ic\_flush\_i}       & in  & 1  & CPU IF stage & Invalidation request pulse \\
  \hline
  \texttt{ic\_stall\_o}       & out & 1  & CPU IF stage & Stall: miss, flush, or error recovery \\
  \hline
  \texttt{ic\_flush\_done\_o} & out & 1  & CPU IF stage & Flush complete pulse \\
  \hline
  \texttt{ic\_err\_o}         & out & 1  & CPU IF stage & AXI4 RRESP error pulse \\
  \hline
  \texttt{ic\_rdata\_o}       & out & 32 & CPU IF stage & Instruction word selected from hit line \\
  \hline
  \texttt{ic\_rvalid\_o}      & out & 1  & CPU IF stage & \texttt{ic\_rdata\_o} valid \\
  \hline
  \multicolumn{5}{|l|}{\textit{I-Tag SRAM port (to \texttt{asICacheBlock})}} \\
  \hline
  \texttt{it\_addr\_o}  & out & 5  & \texttt{asICacheBlock} & Set index from \texttt{ic\_addr\_i[9:5]} \\
  \hline
  \texttt{it\_wdata\_o} & out & 96 & \texttt{asICacheBlock} & Tag word written on fill completion \\
  \hline
  \texttt{it\_we\_o}    & out & 1  & \texttt{asICacheBlock} & Tag write enable \\
  \hline
  \texttt{it\_ce\_o}    & out & 1  & \texttt{asICacheBlock} & Tag chip enable \\
  \hline
  \texttt{it\_rdata\_i} & in  & 96 & \texttt{asICacheBlock} & Tag read data \\
  \hline
  \multicolumn{5}{|l|}{\textit{I-Data SRAM port (to \texttt{asICacheBlock})}} \\
  \hline
  \texttt{id\_addr\_o}  & out & 7  & \texttt{asICacheBlock} & \{index[4:0], way[1:0]\} \\
  \hline
  \texttt{id\_wdata\_o} & out & 64 & \texttt{asICacheBlock} & One AXI4 R-channel beat \\
  \hline
  \texttt{id\_wlane\_o} & out & 2  & \texttt{asICacheBlock} & Beat index (0--3) \\
  \hline
  \texttt{id\_we\_o}    & out & 1  & \texttt{asICacheBlock} & Data write enable \\
  \hline
  \texttt{id\_ce\_o}    & out & 1  & \texttt{asICacheBlock} & Data chip enable \\
  \hline
  \texttt{id\_rdata\_i} & in  & 256 & \texttt{asICacheBlock} & Full 256-bit cache line \\
  \hline
\end{tabularx}
\end{table}

\begin{table}[H]
\caption{\texttt{asICacheCtrl} Port List (2)}
\label{tab:asICacheCtrlPorts2}
\centering
\begin{tabularx}{\textwidth}{|l |l |l |l |X|}
  \hline
  Signal & Dir & Width & Connected to & Description \\
  \hline
  \hline 
  \multicolumn{5}{|l|}{\textit{AXI4 Master Port -- AR and R channels only (to \texttt{asMemArb})}} \\
  \hline
  \texttt{arid\_o}    & out & 4  & \texttt{asMemArb} & Always 4'h1 (I-Cache ID) \\
  \hline
  \texttt{araddr\_o}  & out & 32 & \texttt{asMemArb} & Cache-line-aligned miss address (\texttt{ic\_addr\_i[31:5],5'b0}) \\
  \hline
  \texttt{arlen\_o}   & out & 8  & \texttt{asMemArb} & Always 8'h03 (4 beats) \\
  \hline
  \texttt{arsize\_o}  & out & 3  & \texttt{asMemArb} & Always 3'b011 (8\,B/beat) \\
  \hline
  \texttt{arburst\_o} & out & 2  & \texttt{asMemArb} & Always 2'b01 (INCR) \\
  \hline
  \texttt{arvalid\_o} & out & 1  & \texttt{asMemArb} & AR channel valid; asserted in FILL state \\
  \hline
  \texttt{arready\_i} & in  & 1  & \texttt{asMemArb} & AR channel ready from arbiter \\
  \hline
  \texttt{rid\_i}     & in  & 4  & \texttt{asMemArb} & Checked against 4'h1; mismatch raises \texttt{ic\_err\_o} \\
  \hline
  \texttt{rdata\_i}   & in  & 64 & \texttt{asMemArb} & R-channel beat data \\
  \hline
  \texttt{rresp\_i}   & in  & 2  & \texttt{asMemArb} & 2'b00\,=\,OKAY; error triggers ERR state \\
  \hline
  \texttt{rlast\_i}   & in  & 1  & \texttt{asMemArb} & Last beat indicator \\
  \hline
  \texttt{rvalid\_i}  & in  & 1  & \texttt{asMemArb} & R-channel data valid \\
  \hline
  \texttt{rready\_o}  & out & 1  & \texttt{asMemArb} & Held high throughout FILL \\
  \hline
\end{tabularx}
\end{table}


\paragraph{Internal algorithm.}
The FSM is described in Table~\ref{tab:asICacheFsm}.
The sequential logic is as follows:

\begin{enumerate}
  \item \textbf{IDLE\,$\to$\,LOOKUP:}
    On \texttt{ic\_req\_i}\,=\,1 assert \texttt{it\_ce\_o} with \\
    \texttt{it\_addr\_o}\,=\,\texttt{ic\_addr\_i[9:5]}.
    Advance to LOOKUP.

  \item \textbf{LOOKUP (tag comparison):}
    Read \texttt{it\_rdata\_i[95:0]}.
    For each way $w \in \{0,1,2,3\}$ extract
    \texttt{valid[w]}\,=\,\texttt{it\_rdata\_i}[$23w{+}22$] and
    \texttt{tag[w][21:0]}\,=\,\texttt{it\_rdata\_i}[$23w{+}21\,:\,23w$].
    Compare each \texttt{tag[w]} against \texttt{ic\_addr\_i[31:10]}.
    A \emph{hit} requires \texttt{valid[w]}\,=\,1 and tag match.
    \begin{itemize}
      \item \textbf{Hit:} assert \texttt{id\_ce\_o} with
        \texttt{id\_addr\_o}\,=\,\{index, hit\_way\};
        update PLRU bits in \texttt{it\_wdata\_o} and assert
        \texttt{it\_we\_o}; advance to RESPOND.
      \item \textbf{Miss:} select victim way via PLRU;
        assert \texttt{ic\_stall\_o}; advance to FILL.
    \end{itemize}

  \item \textbf{FILL (AXI4 AR handshake):}
    Drive AR channel with aligned miss address;
    hold \texttt{arvalid\_o}\,=\,1 until \texttt{arready\_i}\,=\,1.
    For each R-channel beat (\texttt{rvalid\_i}\,=\,1):
    write \texttt{rdata\_i} to \texttt{asICacheBlock} via
    \texttt{id\_wdata\_o}, \texttt{id\_wlane\_o}\,=\,beat\_counter[1:0],
    \texttt{id\_we\_o}\,=\,1.
    On \texttt{rlast\_i}: write tag word (\texttt{it\_we\_o}\,=\,1,
    \texttt{it\_wdata\_o} with valid\,=\,1 for the filled way);
    advance to RESPOND.

  \item \textbf{RESPOND:}
    Read \texttt{id\_rdata\_i[255:0]};
    select 32-bit instruction word using \texttt{ic\_addr\_i[4:2]};
    assert \texttt{ic\_rdata\_o} and \texttt{ic\_rvalid\_o};
    de-assert \texttt{ic\_stall\_o}; return to IDLE.

  \item \textbf{FLUSH:}
    Iterate set counter from 0 to 31 (one cycle per set);
    write \texttt{it\_wdata\_o}\,=\,96'h0 (\texttt{it\_we\_o}\,=\,1)
    to clear all valid bits and PLRU bits.
    On counter\,=\,31: pulse \texttt{ic\_flush\_done\_o}; return to IDLE.

  \item \textbf{ERR:}
    Pulse \texttt{ic\_err\_o} for one cycle; de-assert
    \texttt{ic\_stall\_o}; return to IDLE.
\end{enumerate}

\input{\TEXTmems 04_memsyst_icache_fsm}

% -------------------------------------------------------------------------
\subsubsection{\texttt{asDCacheCtrl} -- Data Cache Controller}
\label{subsubsec:asDCacheCtrl}

\paragraph{Purpose.}
Implements the D-Cache FSM (Table~\ref{tab:asDCacheFsm}).
Serves CPU load/store requests via \texttt{as\_dcache\_if};
routes Flash-region loads through the \texttt{asDCacheBlock} SRAM and
the AXI4 miss path; routes Scratchpad-region loads and stores directly
to \texttt{asSpBlock} (bypass path).
No write path to Flash exists.

\paragraph{Port list} of the D-cache controller:

\begin{table}[H]
\caption{\texttt{asDCacheCtrl} Port List (1)}
\label{tab:asDCacheCtrlPorts}
\centering
\begin{tabularx}{\textwidth}{|l |l |l |l |X|}
  \hline
  Signal & Dir & Width & Connected to & Description \\
  \hline
  \hline
  \texttt{clk\_i}  & in & 1 & \texttt{asMemTop} & System clock \\
  \hline
  \texttt{rst\_i}  & in & 1 & \texttt{asMemTop} & Synchronous reset \\
  \hline
  \multicolumn{5}{|l|}{\textit{CPU Data Bus (\texttt{as\_dcache\_if.cache} modport)}} \\
  \hline
  \texttt{dc\_addr\_i}        & in  & 64 & CPU MEM stage & Load/store address \\
  \hline
  \texttt{dc\_req\_i}         & in  & 1  & CPU MEM stage & Request valid \\
  \hline
  \texttt{dc\_wr\_i}          & in  & 1  & CPU MEM stage & 1\,=\,store, 0\,=\,load \\
  \hline
  \texttt{dc\_size\_i}        & in  & 3  & CPU MEM stage & Transfer size \\
  \hline
  \texttt{dc\_wdata\_i}       & in  & 64 & CPU MEM stage & Store data \\
  \hline
  \texttt{dc\_wstrb\_i}       & in  & 8  & CPU MEM stage & Byte enables \\
  \hline
  \texttt{dc\_flush\_i}       & in  & 1  & CPU MEM stage & Invalidation pulse \\
  \hline
  \texttt{dc\_stall\_o}       & out & 1  & CPU MEM stage & Stall signal \\
  \hline
  \texttt{dc\_flush\_done\_o} & out & 1  & CPU MEM stage & Flush complete pulse \\
  \hline
  \texttt{dc\_err\_o}         & out & 1  & CPU MEM stage & AXI4 error pulse \\
  \hline
  \texttt{dc\_rdata\_o}       & out & 64 & CPU MEM stage & Load data \\
  \hline
  \texttt{dc\_rvalid\_o}      & out & 1  & CPU MEM stage & Load data valid \\
  \hline
  \multicolumn{5}{|l|}{\textit{D-Tag SRAM port (to \texttt{asDCacheBlock})}} \\
  \hline
  \texttt{dt\_addr\_o}  & out & 5   & \texttt{asDCacheBlock} & Set index \\
  \hline
  \texttt{dt\_wdata\_o} & out & 96  & \texttt{asDCacheBlock} & Tag write data \\
  \hline
  \texttt{dt\_we\_o}    & out & 1   & \texttt{asDCacheBlock} & Tag write enable \\
  \hline
  \texttt{dt\_ce\_o}    & out & 1   & \texttt{asDCacheBlock} & Tag chip enable \\
  \hline
  \texttt{dt\_rdata\_i} & in  & 96  & \texttt{asDCacheBlock} & Tag read data \\
  \hline
  \multicolumn{5}{|l|}{\textit{D-Data SRAM port (to \texttt{asDCacheBlock})}} \\
  \hline
  \texttt{dd\_addr\_o}  & out & 7   & \texttt{asDCacheBlock} & \{index[4:0], way[1:0]\} \\
  \hline
  \texttt{dd\_wdata\_o} & out & 64  & \texttt{asDCacheBlock} & AXI4 R-channel beat \\
  \hline
  \texttt{dd\_wlane\_o} & out & 2   & \texttt{asDCacheBlock} & Beat index (0--3) \\
  \hline
  \texttt{dd\_we\_o}    & out & 1   & \texttt{asDCacheBlock} & Data write enable \\
  \hline
  \texttt{dd\_ce\_o}    & out & 1   & \texttt{asDCacheBlock} & Data chip enable \\
  \hline
  \texttt{dd\_rdata\_i} & in  & 256 & \texttt{asDCacheBlock} & Full cache line \\
  \hline
\end{tabularx}
\end{table}

\begin{table}[H]
\caption{\texttt{asDCacheCtrl} Port List (2)}
\label{tab:asDCacheCtrlPorts2}
\centering
\begin{tabularx}{\textwidth}{|l |l |l |l |X|}
  \hline
  Signal & Dir & Width & Connected to & Description \\
  \hline
  \hline
  \multicolumn{5}{|l|}{\textit{Scratchpad SRAM port (to \texttt{asSpBlock})}} \\
  \hline
  \texttt{sp\_addr\_o}  & out & 10 & \texttt{asSpBlock} & Word address \texttt{dc\_addr\_i[12:3]} \\
  \hline
  \texttt{sp\_wdata\_o} & out & 64 & \texttt{asSpBlock} & Store data from \texttt{dc\_wdata\_i} \\
  \hline
  \texttt{sp\_wstrb\_o} & out & 8  & \texttt{asSpBlock} & Byte enables from \texttt{dc\_wstrb\_i} \\
  \hline
  \texttt{sp\_we\_o}    & out & 1  & \texttt{asSpBlock} & Write enable: high on store to Scratchpad region \\
  \hline
  \texttt{sp\_ce\_o}    & out & 1  & \texttt{asSpBlock} & Chip enable: high on any Scratchpad access \\
  \hline
  \texttt{sp\_rdata\_i} & in  & 64 & \texttt{asSpBlock} & Load data from Scratchpad \\
  \hline
  \multicolumn{5}{|l|}{\textit{AXI4 Master Port -- AR and R channels only (to \texttt{asMemArb})}} \\
  \hline
  \texttt{arid\_o}    & out & 4  & \texttt{asMemArb} & Always 4'h2 (D-Cache ID) \\
  \hline
  \texttt{araddr\_o}  & out & 32 & \texttt{asMemArb} & Cache-line-aligned miss address \\
  \hline
  \texttt{arlen\_o}   & out & 8  & \texttt{asMemArb} & Always 8'h03 \\
  \hline
  \texttt{arsize\_o}  & out & 3  & \texttt{asMemArb} & Always 3'b011 \\
  \hline
  \texttt{arburst\_o} & out & 2  & \texttt{asMemArb} & Always 2'b01 (INCR) \\
  \hline
  \texttt{arvalid\_o} & out & 1  & \texttt{asMemArb} & AR valid; asserted in FILL \\
  \hline
  \texttt{arready\_i} & in  & 1  & \texttt{asMemArb} & AR ready \\
  \hline
  \texttt{rid\_i}     & in  & 4  & \texttt{asMemArb} & Checked against 4'h2 \\
  \hline
  \texttt{rdata\_i}   & in  & 64 & \texttt{asMemArb} & R-channel beat \\
  \hline
  \texttt{rresp\_i}   & in  & 2  & \texttt{asMemArb} & Response; OKAY\,=\,2'b00 \\
  \hline
  \texttt{rlast\_i}   & in  & 1  & \texttt{asMemArb} & Last beat \\
  \hline
  \texttt{rvalid\_i}  & in  & 1  & \texttt{asMemArb} & R-channel valid \\
  \hline
  \texttt{rready\_o}  & out & 1  & \texttt{asMemArb} & Held high throughout FILL \\
  \hline
\end{tabularx}
\end{table}



\paragraph{Internal algorithm.}
The FSM is described in Table~\ref{tab:asDCacheFsm}.
Address decoding in the DECODE state is the critical dispatch:

\begin{enumerate}
  \item \textbf{DECODE -- address region check:}
    \texttt{dc\_addr\_i[31:13]} is compared against
    \texttt{SP\_BASE[31:13]}.
    \begin{itemize}
      \item Match $\Rightarrow$ Scratchpad region: advance to BYPASS.
      \item No match $\Rightarrow$ Flash region: assert \texttt{dt\_ce\_o}
        and advance to LOOKUP.
    \end{itemize}

  \item \textbf{BYPASS (Scratchpad access):}
    For a \textbf{store}: drive \texttt{sp\_addr\_o},
    \texttt{sp\_wdata\_o}, \texttt{sp\_wstrb\_o},
    \texttt{sp\_we\_o}\,=\,1, \texttt{sp\_ce\_o}\,=\,1.
    De-assert \texttt{dc\_stall\_o} after one cycle; return to IDLE.
    For a \textbf{load}: drive \texttt{sp\_addr\_o},
    \texttt{sp\_we\_o}\,=\,0, \texttt{sp\_ce\_o}\,=\,1.
    Latch \texttt{sp\_rdata\_i} into \texttt{dc\_rdata\_o} in the
    following cycle; assert \texttt{dc\_rvalid\_o}; return to IDLE.

  \item \textbf{LOOKUP -- Flash region, tag comparison:}
    Same tag-compare logic as \texttt{asICacheCtrl}.
    \begin{itemize}
      \item \textbf{Hit:} read \texttt{dd\_rdata\_i[255:0]};
        extract 64-bit double word at
        \texttt{dc\_addr\_i[4:3]}\,$\times$\,64\,bit;
        sign/zero-extend according to \texttt{dc\_size\_i};
        assert \texttt{dc\_rdata\_o}, \texttt{dc\_rvalid\_o};
        update PLRU; return to IDLE.
      \item \textbf{Miss:} assert \texttt{dc\_stall\_o}; advance to FILL.
    \end{itemize}
    \textbf{Note:} Stores to the Flash region are not permitted.
    If \texttt{dc\_wr\_i}\,=\,1 and address is in Flash region the
    controller asserts \texttt{dc\_err\_o} and returns to IDLE without
    any SRAM or AXI4 access.

  \item \textbf{FILL:}
    Identical to \texttt{asICacheCtrl} FILL, using D-Cache AXI4 ID
    4'h2 and \texttt{asDCacheBlock} SRAM ports.

  \item \textbf{RESPOND, FLUSH, ERR:}
    Identical in structure to \texttt{asICacheCtrl}; see
    Section~\ref{subsubsec:asICacheCtrl}.
\end{enumerate}

\input{\TEXTmems 04_memsyst_dcache_fsm}

% -------------------------------------------------------------------------
\subsubsection{\texttt{asMemArb} -- Memory Arbiter}
\label{subsubsec:asMemArb}

\paragraph{Purpose.}
Fixed-priority combinatorial AXI4 multiplexer.
Serialises the AR/R (and potentially AW/W/B) channels of the I-Cache
and D-Cache controllers onto the single AXI4 port of the QSPI
controller.
D-Cache has priority over I-Cache (Table~\ref{tab:asArbPriority}).

\paragraph{Port list} of the memory arbiter:

\begin{table}[H]
\caption{\texttt{asMemArb} Port List (1)}
\label{tab:asMemArbPorts}
\centering
\begin{tabularx}{\textwidth}{|X |l |l |X |X|}
  \hline
  Signal & Dir & Width & Connected to & Description \\
  \hline
  \hline
  \texttt{clk\_i}  & in & 1 & \texttt{asMemTop} & System clock (used only for grant register) \\
  \hline
  \texttt{rst\_i}  & in & 1 & \texttt{asMemTop} & Synchronous reset; clears grant register \\
  \hline
  \multicolumn{5}{|l|}{\textit{I-Cache AXI4 Slave Port (AR + R channels)}} \\
  \hline
  \texttt{ic\_arvalid\_i} & in  & 1  & \texttt{asICacheCtrl} & I-Cache AR valid \\
  \hline
  \texttt{ic\_araddr\_i}  & in  & 32 & \texttt{asICacheCtrl} & I-Cache AR address \\
  \hline
  \texttt{ic\_arid\_i}    & in  & 4  & \texttt{asICacheCtrl} & Always 4'h1 \\
  \hline
  \texttt{ic\_arlen\_i}   & in  & 8  & \texttt{asICacheCtrl} & Always 8'h03 \\
  \hline
  \texttt{ic\_arsize\_i}  & in  & 3  & \texttt{asICacheCtrl} & Always 3'b011 \\
  \hline
  \texttt{ic\_arburst\_i} & in  & 2  & \texttt{asICacheCtrl} & Always 2'b01 \\
  \hline
  \texttt{ic\_arready\_o} & out & 1  & \texttt{asICacheCtrl} & AR ready back to I-Cache \\
  \hline
  \texttt{ic\_rvalid\_o}  & out & 1  & \texttt{asICacheCtrl} & R-channel valid routed to I-Cache \\
  \hline
  \texttt{ic\_rdata\_o}   & out & 64 & \texttt{asICacheCtrl} & R-channel data routed to I-Cache \\
  \hline
  \texttt{ic\_rresp\_o}   & out & 2  & \texttt{asICacheCtrl} & R-channel response to I-Cache \\
  \hline
  \texttt{ic\_rlast\_o}   & out & 1  & \texttt{asICacheCtrl} & R-channel last beat to I-Cache \\
  \hline
  \texttt{ic\_rid\_o}     & out & 4  & \texttt{asICacheCtrl} & R-channel ID to I-Cache \\
  \hline
  \texttt{ic\_rready\_i}  & in  & 1  & \texttt{asICacheCtrl} & R-channel ready from I-Cache \\
  \hline
  \multicolumn{5}{|l|}{\textit{D-Cache AXI4 Slave Port (AR + R channels; identical structure)}} \\
  \hline
  \texttt{dc\_ar*\_i / dc\_r*\_o / dc\_rready\_i} & in/out & -- & \texttt{asDCacheCtrl} & Same signals as I-Cache port; prefix \texttt{dc\_} \\
  \hline
\end{tabularx}
\end{table}

\begin{table}[H]
\caption{\texttt{asMemArb} Port List (2)}
\label{tab:asMemArbPorts2}
\centering
\begin{tabularx}{\textwidth}{|X |l |l |X |X|}
  \hline
  Signal & Dir & Width & Connected to & Description \\
  \hline
  \hline
  \multicolumn{5}{|l|}{\textit{QSPI Controller AXI4 Master Port}} \\
  \hline
  \texttt{qspi\_arvalid\_o} & out & 1  & \texttt{asQspiTop} & AR valid to QSPI \\
  \hline
  \texttt{qspi\_araddr\_o}  & out & 32 & \texttt{asQspiTop} & AR address to QSPI \\
  \hline
  \texttt{qspi\_arid\_o}    & out & 4  & \texttt{asQspiTop} & AR ID forwarded from winning master \\
  \hline
  \texttt{qspi\_arlen\_o}   & out & 8  & \texttt{asQspiTop} & Always 8'h03 \\
  \hline
  \texttt{qspi\_arsize\_o}  & out & 3  & \texttt{asQspiTop} & Always 3'b011 \\
  \hline
  \texttt{qspi\_arburst\_o} & out & 2  & \texttt{asQspiTop} & Always 2'b01 \\
  \hline
  \texttt{qspi\_arready\_i} & in  & 1  & \texttt{asQspiTop} & AR ready from QSPI \\
  \hline
  \texttt{qspi\_rvalid\_i}  & in  & 1  & \texttt{asQspiTop} & R-channel valid from QSPI \\
  \hline
  \texttt{qspi\_rdata\_i}   & in  & 64 & \texttt{asQspiTop} & R-channel data from QSPI \\
  \hline
  \texttt{qspi\_rresp\_i}   & in  & 2  & \texttt{asQspiTop} & R-channel response from QSPI \\
  \hline
  \texttt{qspi\_rlast\_i}   & in  & 1  & \texttt{asQspiTop} & R-channel last beat from QSPI \\
  \hline
  \texttt{qspi\_rid\_i}     & in  & 4  & \texttt{asQspiTop} & R-channel ID from QSPI \\
  \hline
  \texttt{qspi\_rready\_o}  & out & 1  & \texttt{asQspiTop} & R-channel ready to QSPI (OR of both master readys) \\
  \hline
\end{tabularx}
\end{table}

\paragraph{Internal algorithm.}
The arbiter maintains a one-bit \textbf{grant register}
(\texttt{grant\_r}: 0\,=\,I-Cache, 1\,=\,D-Cache) and a
\textbf{busy flag} (\texttt{busy\_r}) that locks the grant for the
duration of an active burst.

\begin{enumerate}
  \item \textbf{IDLE} (\texttt{busy\_r}\,=\,0):
    \begin{itemize}
      \item If \texttt{dc\_arvalid\_i}\,=\,1:
        set \texttt{grant\_r}\,=\,1, \texttt{busy\_r}\,=\,1;
        forward D-Cache AR signals to QSPI master port.
      \item Else if \texttt{ic\_arvalid\_i}\,=\,1:
        set \texttt{grant\_r}\,=\,0, \texttt{busy\_r}\,=\,1;
        forward I-Cache AR signals to QSPI master port.
    \end{itemize}

  \item \textbf{BUSY} (\texttt{busy\_r}\,=\,1):
    All QSPI AR/R signals are forwarded to the granted master only.
    \texttt{arready\_o} of the non-granted master is held low (bus
    not available).
    \texttt{busy\_r} is cleared when \texttt{qspi\_rlast\_i}\,=\,1
    and \texttt{qspi\_rvalid\_i}\,=\,1
    (last beat of the burst received).

  \item \textbf{R-channel routing:}
    \texttt{qspi\_rid\_i} identifies the receiving master (4'h1\,=\,I-Cache,
    4'h2\,=\,D-Cache); R-channel signals are forwarded to the
    matching master regardless of \texttt{grant\_r}, ensuring correct
    response routing even during back-to-back grants.
\end{enumerate}

The arbiter is purely combinatorial on the data path; only
\texttt{grant\_r} and \texttt{busy\_r} are registered.

% -------------------------------------------------------------------------
\subsubsection{\texttt{asSpramModel} -- Simulation SRAM Behavioural Model}
\label{subsubsec:asSpramModel}

\paragraph{Purpose.}
Provides a synthesisable-free behavioural model of the X-Fab XO035
SPRAM macro for use in RTL simulation (Vivado \texttt{xsim}, ModelSim,
VCS).
It is \textbf{not} used in synthesis; the three SRAM wrapper modules
select between the real SPRAM macro and \texttt{asSpramModel} via a
conditional compile guard:

\begin{verbatim}
`ifdef SIMULATION
  asSpramModel #(.WIDTH(W), .DEPTH(D)) u_sram (...);
`else
  SPRAM_XO035   #(.WIDTH(W), .DEPTH(D)) u_sram (...);
`endif
\end{verbatim}

\paragraph{Port list.}

\begin{table}[H]
\caption{\texttt{asSpramModel} Port List (matches XO035 SPRAM interface)}
\label{tab:asSpramModelPorts}
\centering
\begin{tabularx}{\textwidth}{|l |l |l |X|}
  \hline
  Signal & Dir & Width & Description \\
  \hline
  \hline
  \texttt{CLK}  & in  & 1     & Rising-edge clock \\
  \hline
  \texttt{CEN}  & in  & 1     & Chip enable, active-low; high\,=\,standby (no access) \\
  \hline
  \texttt{WEN}  & in  & 1     & Write enable, active-low; high\,=\,read, low\,=\,write \\
  \hline
  \texttt{BWEN} & in  & WIDTH & Bit write enable, active-low; 0\,=\,write bit, 1\,=\,mask \\
  \hline
  \texttt{A}    & in  & $\lceil\log_2(\text{DEPTH})\rceil$ & Address \\
  \hline
  \texttt{D}    & in  & WIDTH & Write data \\
  \hline
  \texttt{Q}    & out & WIDTH & Read data; updated one clock after access \\
  \hline
\end{tabularx}
\end{table}

\paragraph{Internal algorithm.}
\begin{verbatim}
parameter int WIDTH = 96;
parameter int DEPTH = 32;
logic [WIDTH-1:0] mem [0:DEPTH-1];

always_ff @(posedge CLK) begin
  if (!CEN) begin
    if (!WEN)                          // write
      for (int b = 0; b < WIDTH; b++)
        if (!BWEN[b]) mem[A][b] <= D[b];
    Q <= mem[A];                       // read (write-first)
  end
end
\end{verbatim}

The model uses \emph{write-first} semantics: on a write cycle
\texttt{Q} returns the newly written data, matching the XO035 SPRAM
read-during-write behaviour.
The \texttt{initial} block zero-initialises \texttt{mem} to
model power-on state.
%%%%%%%%%%%%%%%%%%%%%%%%%%%%%%%%%%%%%%%%%%%%%%%%%%%%%%%%%%%






%%%%%%%%%%%%%%%%%%%%%%%%%%%%%%%%%
% Service Requests of the MemSys
%%%%%%%%%%%%%%%%%%%%%%%%%%%%%%%%%
\subsection{Service Requests of the MemSys}
No


%%%%%%%%%%%%%%%%%%%%%%%%%%%%%%%%%
% The MemSys Peripheral Kernel
%%%%%%%%%%%%%%%%%%%%%%%%%%%%%%%%%
\subsection{The MemSys Peripheral Kernel}
\label{subsec:memsysfunc02}
\paragraph{Block Architecture: } 



%%%%%%%%%%%%%%%%%%%%%%%%%%%%%%%%%%%%%%%%%%
% All Registers
%%%%%%%%%%%%%%%%%%%%%%%%%%%%%%%%%%%%%%%%%%
\subsection{Registers}
No registers.

